% Stage 8 restored fragment — 2.8.3 Reciprocity in radiation and reception (archive lines 10909-11155)
\subsection{Reciprocity in radiation and reception}
\label{sec:ch283}

A standard-gain horn on a compact-range turntable fixes the laboratory scene for `Reciprocity in radiation and reception'': frequency, probe standoff, and wattmeter closure on $\Gamma_R$ precede the formal displays below \cite{hansen1988}. Maxwell's retarded export to $\Gamma_R$ remains the governing constraint \cite{maxwell1873}.
Two-port antenna measurements are the everyday face of reciprocity. A horn coupled to a compact
range, a patch fed from a coax port, or a dual-polarized reflector feed all reduce---after
calibration---to the same question the vector network analyzer answers: does the complex coupling
from port~$A$ to port~$B$ equal the reversed coupling when cables, standards, and probe
normalization are exchanged? Subsection~\ref{sec:ch283} gives the Chapter~\ref{ch:02} definitive
account of that interchange in operator form: transmit--receive operator symmetry, two-port
coupling bilinears, and measurement-qualified reciprocity on
\(\operatorname{ran}\mathcal{R}_\omega\), developing Eq.
Eq.~\eqref{eq:ch2.transmit-receive-swap} \cite{stratton1941}. Observation
composition Eq.~\eqref{eq:ch2.reconstruction-observation-map}, finite-rank readouts
Eq.~\eqref{eq:ch2.finite-rank-observation}, rank cap
Eq.~\eqref{eq:ch2.reconstruction-rank-cap}, and null-space laws
Eqs.~\eqref{eq:ch2.reconstruction-null-space}--\eqref{eq:ch2.nr-superposition-law} are
presupposed; inverse-source regularization is not re-developed \cite{stratton1941}.
Plane-wave angular-spectrum filters belong to Section~\ref{sec:ch29}; explicit port circuit
models are not tabulated \cite{stratton1941}.
Assumptions. Fix \(\omega>0\), reciprocal media Eq.~\eqref{eq:ch2.reciprocal-medium-class}, passive
boundary Eq.~\eqref{eq:ch2.passive-impedance-class}, outgoing radiation condition, and admissible sources
\(\Jimp^{\mathrm{(tx)}},\Jimp^{\mathrm{(rx)}}\in\mathcal{J}_\Omega\) on ports \(A\) and \(B\)
\cite{stratton1941}. Each port carries a finite-rank observation functional
\(\mathcal{O}_{\mathrm{meas}}^{(P)}:\mathcal{F}_{\mathrm{rad}}\to\mathbb{C}\) or
\(\mathbb{C}^{N_P}\) of the form Eq.~\eqref{eq:ch2.finite-rank-observation}
\cite{stratton1941}. regime qualification Eq.~\eqref{eq:ch2.regime-operator-screen} and transport
qualification Eq.~\eqref{eq:ch2.transport-qualified-radiation} govern when port outputs may be
reported as radiative transport \cite{harrington2001}. The same \(\mathcal{B}_\Gamma\), outgoing
branch, and material data are used for both excitations \cite{stratton1941}.
For port \(P\in\{A,B\}\), define the \emph{port radiation--reception operator}

\begin{equation}
\label{eq:ch2.port-radiation-reception-operator}
\mathcal{M}_P
=
\mathcal{O}_{\mathrm{meas}}^{(P)}\circ\Lambda_{\mathrm{rad}},
\qquad
\Lambda_{\mathrm{rad}}=\mathcal{O}_{\mathrm{far}}\circ\mathcal{R}_\omega,
\end{equation}
consistent with Eq.~\eqref{eq:ch2.reconstruction-observation-map}
\cite{stratton1941}. When \(\mathcal{O}_{\mathrm{meas}}^{(P)}\) is scalar,
\(\mathcal{M}_P[\Jimp]\in\mathbb{C}\) is the complex port response induced by impressed
\(\Jimp\); when vector-valued, \(\mathcal{M}_P[\Jimp]\in\mathbb{C}^{N_P}\) lists channel
readouts \cite{stratton1941}. The \emph{transmit map} at port \(B\) due to excitation at
port \(A\) is the composition \(\Jimp^{\mathrm{(tx)}}\mapsto
\mathcal{M}_B[\Jimp^{\mathrm{(tx)}}]\); the \emph{receive map} at port \(A\) due to excitation
at port \(B\) is \(\Jimp^{\mathrm{(rx)}}\mapsto\mathcal{M}_A[\Jimp^{\mathrm{(rx)}}]\)
\cite{harrington2001}. Algebraically, both are finite-rank linear functionals of
\(\mathcal{R}_\omega[\Jimp]\); physically, reciprocity constrains their interchange when
hardware and media are reciprocal \cite{stratton1941}.
Under Lorentz reciprocity for the solution map (Subsection~\ref{sec:ch281}) and self-adjoint port
functionals \(\mathcal{O}_{\mathrm{meas}}^{(P)}\) on \(\mathcal{F}_{\mathrm{rad}}\), the
complex port response at \(B\) induced by excitation at \(A\) equals the response at \(A\) when
the impressed source and receive probe are exchanged on reciprocal hardware:
\begin{equation}
\label{eq:ch2.transmit-receive-reciprocity}
\mathcal{M}_B\big[\Jimp^{\mathrm{(tx)}}\big]
=
\mathcal{M}_A\big[\Jimp^{\mathrm{(rx)}}\big]
\quad\text{when}\ 
\Jimp^{\mathrm{(tx)}}\ \text{and}\ \Jimp^{\mathrm{(rx)}}\ \text{are reciprocal port excitations},
\end{equation}
developing Eq.~\eqref{eq:ch2.transmit-receive-swap}
\cite{stratton1941}. Equation~\eqref{eq:ch2.transmit-receive-reciprocity}
is the operator reading of the VNA experiment: excite \(A\), read \(B\); swap roles on reciprocal
media and obtain the same complex port coupling on observable coordinates. It states that the
observable induced at \(B\) by the transmit source at \(A\) equals the observable at \(A\) when
the roles of impressed source and probe are exchanged on reciprocal media
\cite{stratton1941}. It is not a claim that \(\mathcal{M}_B=\mathcal{M}_A\) as
operators on \(\mathcal{J}_\Omega\): ports may have distinct apertures and channel counts
\cite{harrington2001}. It \emph{is} a claim that the \emph{bilinear} port coupling is symmetric
when qualified by Eq.~\eqref{eq:ch2.radiative-coupling-bilinear} on
\(\operatorname{ran}\mathcal{R}_\omega\) \cite{stratton1941}.
\begin{figure}[t]
\centering
\ChOneFigBlock{%
\begin{tikzpicture}[ch1fig, node distance=5mm and 7mm]
 \node[ch1box] (tx) {Port \(A\)\\$\Jimp^{\mathrm{(tx)}}$};
 \node[ch1box, right=7mm of tx] (r) {$\mathcal{R}_\omega$};
 \node[ch1box, right=7mm of r] (ob) {$\mathcal{O}_{\mathrm{meas}}^{(B)}$};
 \node[ch1box, right=7mm of ob] (rxout) {Rx at \(B\)};
 \node[ch1box, below=5mm of tx] (rx) {Port \(B\)\\$\Jimp^{\mathrm{(rx)}}$};
 \node[ch1box, below=5mm of ob] (ob2) {$\mathcal{O}_{\mathrm{meas}}^{(A)}$};
 \node[ch1box, below=5mm of rxout] (txout) {Rx at \(A\)};
 \draw[ch1flow] (tx) -- (r) -- (ob) -- (rxout);
 \draw[ch1flow] (rx) -- (r) -- (ob2) -- (txout);
 \draw[ch1flow, dashed] (rxout) to[bend left=20] node[right,font=\scriptsize] {swap} (txout);
\end{tikzpicture}%
}
\caption{Subsection~\ref{sec:ch283} transmit--receive measurement sequence: excitation at one port and
readout at the other interchange under Eq.~\eqref{eq:ch2.transmit-receive-reciprocity} on
reciprocal media.}
\label{fig:ch2.transmit-receive-measurement chain}
\end{figure}
Figure interpretation. Transmit port $\Jimp^{\mathrm{(tx)}}$ and receive probe $\Jimp^{\mathrm{(rx)}}$ compose through the same $\mathcal{R}_\omega$ and observation $\mathcal{O}_{\mathrm{meas}}$. Swap symmetry Eq.~\eqref{eq:ch2.transmit-receive-swap} holds on observable coordinates after projection onto $\operatorname{ran}\mathcal{R}_\omega$; additions in $\ker\mathcal{R}_\omega$ leave port coupling unchanged. The schematic separates reciprocity-qualified S-parameters from raw near-field cross terms that may violate apparent symmetry while far-zone reciprocity remains valid \cite{stratton1941}. Figure~\ref{fig:ch2.transmit-receive-measurement chain} displays the two compositions constrained by
Eq.~\eqref{eq:ch2.transmit-receive-reciprocity}. Both paths pass through the same
\(\mathcal{R}_\omega\); port asymmetry enters only through \(\mathcal{O}_{\mathrm{meas}}^{(A)}\),
\(\mathcal{O}_{\mathrm{meas}}^{(B)}\) \cite{stratton1941}.
\phantomsection\label{thm:ch2.radiation-reception-reciprocity}
Let media satisfy Eq.~\eqref{eq:ch2.reciprocal-medium-class} with passive
Eq.~\eqref{eq:ch2.passive-impedance-class} completion, and let port operators
\(\mathcal{M}_P=\mathcal{O}_{\mathrm{meas}}^{(P)}\circ\Lambda_{\mathrm{rad}}\) be defined by
Eq.~\eqref{eq:ch2.port-radiation-reception-operator} with self-adjoint
\(\mathcal{O}_{\mathrm{meas}}^{(P)}\) on \(\mathcal{F}_{\mathrm{rad}}\). For admissible port
sources \(\Jimp^{\mathrm{(tx)}}\), \(\Jimp^{\mathrm{(rx)}}\):
\begin{enumerate}
\item[(i)] Scalar reciprocity Eq.~\eqref{eq:ch2.transmit-receive-reciprocity} holds when port
excitations are paired through Lorentz reciprocity for the solution map (Subsection~\ref{sec:ch281}).
\item[(ii)] Bilinear port coupling
\begin{equation}
\label{eq:ch2.port-coupling-bilinear}
\mathfrak{C}_{AB}
=
\langle \mathcal{O}_{\mathrm{meas}}^{(B)}[\mathcal{R}_\omega[\Jimp^{\mathrm{(tx)}}]],\,
\mathcal{O}_{\mathrm{meas}}^{(A)}[\mathcal{R}_\omega[\Jimp^{\mathrm{(rx)}}]]\rangle
\end{equation}
is symmetric under interchange of port labels when inner products are those induced by
Eq.~\eqref{eq:ch2.field-energy-norm} on \(\mathcal{F}_{\mathrm{rad}}\).
\item[(iii)] Two-port impedance matrix \(\mathbf{Z}\) with entries
\(Z_{PQ}=\mathcal{M}_P[\Jimp^{(Q)}]\) on matched reciprocal ports satisfies
\begin{equation}
\label{eq:ch2.reciprocal-impedance-matrix}
Z_{PQ}=Z_{QP}.
\end{equation}
\item[(iv)] Adding \(\Jimp^{\mathrm{(NR)}}\in\ker\mathcal{R}_\omega\) to either port excitation
leaves Eqs.~\eqref{eq:ch2.transmit-receive-reciprocity}--\eqref{eq:ch2.reciprocal-impedance-matrix}
unchanged on observable coordinates.
\end{enumerate}
Part (i) applies Lorentz reciprocity for the solution map (Subsection~\ref{sec:ch281})(iv) to
\(\mathcal{R}_\omega[\Jimp^{\mathrm{(tx)}}]\), \(\mathcal{R}_\omega[\Jimp^{\mathrm{(rx)}}]\),
then self-adjointness of \(\mathcal{O}_{\mathrm{meas}}^{(P)}\)
\cite{stratton1941}. Part (ii) is
Eq.~\eqref{eq:ch2.coupling-bilinear-symmetry} with observation weights inserted; symmetry of
\(\mathfrak{C}_\omega\) from orthogonality and coupling of radiative modes (Subsection~\ref{sec:ch282})(iii)
\cite{stratton1941}. Part (iii) specializes (ii) to matched port basis vectors
\cite{harrington2001}. Part (iv) uses
\(\mathcal{R}_\omega[\Jimp^{\mathrm{(NR)}}]=\mathbf{0}\) and
null-space invariance of radiative reciprocity (Subsection~\ref{sec:ch281}) \cite{stratton1941}.
Mathematical role of Eq.~\eqref{eq:ch2.port-coupling-bilinear}. Equation~\eqref{eq:ch2.port-coupling-bilinear} packages mutual coupling as a bilinear form on
\(\operatorname{ran}\mathcal{R}_\omega\) after port projection \cite{stratton1941}. It reduces
to Eq.~\eqref{eq:ch2.radiative-coupling-bilinear} when
\(\mathcal{O}_{\mathrm{meas}}^{(A)}=\mathcal{O}_{\mathrm{meas}}^{(B)}=\mathrm{id}\) on the
observed subspace \cite{stratton1941}. Physically, \(\mathfrak{C}_{AB}\) is the
operator form of mutual impedance or scattering amplitude between declared ports when regime qualification
Eq.~\eqref{eq:ch2.regime-component-criterion} is satisfied \cite{harrington2001}. Finite
rank Eq.~\eqref{eq:ch2.reconstruction-rank-cap} limits the number of independent entries
\(\mathfrak{C}_{AB}\) can resolve \cite{stratton1941}.
On orthonormal modes Eq.~\eqref{eq:ch2.radiation-modal-orthogonality}, port responses expand as
\begin{equation}
\label{eq:ch2.port-modal-response-expansion}
\mathcal{M}_P[\Jimp]
=
\sum_{n\in\mathcal{N}_{\mathrm{rad}}}
w_{P,n}\,\lambda_n\,c_n[\Jimp],
\end{equation}
with port weights \(w_{P,n}=\langle \mathbf{f}_P,\mathbf{u}_n\rangle_{\mathcal{U}}\) inherited
from Eq.~\eqref{eq:ch2.finite-rank-observation} and amplitudes \(c_n\) from
Eq.~\eqref{eq:ch2.source-modal-amplitude} \cite{stratton1941}. Mutual coupling
between ports \(A\) and \(B\) becomes
\begin{equation}
\label{eq:ch2.port-modal-coupling-sum}
\mathfrak{C}_{AB}
=
\sum_{n\in\mathcal{N}_{\mathrm{rad}}}
|\lambda_n|^{2}\,w_{A,n}^{*}\,w_{B,n}\,
c_n[\Jimp^{\mathrm{(tx)}}]^{*}\,c_n[\Jimp^{\mathrm{(rx)}}],
\end{equation}
using Eq.~\eqref{eq:ch2.modal-coupling-sum-diagonal} \cite{stratton1941}. Equation
\eqref{eq:ch2.port-modal-coupling-sum} shows that reciprocity symmetry is diagonal in a mode
basis but may be dense in fixed laboratory coordinates when degenerate multiplets from
Subsection~\ref{sec:ch262} are expressed in non-eigenvector bases
\cite{harrington2001}.
\phantomsection\label{prop:ch2.measurement-reciprocity}
For full rank-$N$ instrument readout \(\mathcal{M}_\omega=\mathcal{O}_{\mathrm{meas}}\circ\Lambda_{\mathrm{rad}}\)
Eq.~\eqref{eq:ch2.reconstruction-observation-map}, reciprocity implies
\(\mathcal{M}_\omega[\Jimp^{\mathrm{(tx)}}]\) and \(\mathcal{M}_\omega[\Jimp^{\mathrm{(rx)}}]\)
obey the same swap law Eq.~\eqref{eq:ch2.transmit-receive-reciprocity} on
\(\operatorname{ran}\mathcal{O}_{\mathrm{meas}}\) \cite{stratton1941}. Sources
in \(\ker\mathcal{M}_\omega\) leave all port readouts unchanged
Eq.~\eqref{eq:ch2.reconstruction-null-space}; reciprocity does not resolve those equivalence
classes \cite{stratton1941}. Rank cap Eq.~\eqref{eq:ch2.reconstruction-rank-cap} bounds
independent reciprocal channels \cite{harrington2001}.
Apply reciprocity in radiation and reception (Subsection~\ref{sec:ch283})(i) with
\(\mathcal{O}_{\mathrm{meas}}^{(A)}=\mathcal{O}_{\mathrm{meas}}^{(B)}=\mathcal{O}_{\mathrm{meas}}\)
\cite{stratton1941}. Null-space invariance is
measurement null-space extension (Subsection~\ref{sec:ch271}) \cite{stratton1941}. \phantomsection\label{prop:ch2.accessible-sector-reciprocity}
Let accessible modal sector Eq.~\eqref{eq:ch2.accessible-modal-sector} restrict far-zone content
to \(\mathcal{K}_{\mathrm{acc}}\). Then Eqs.~\eqref{eq:ch2.transmit-receive-reciprocity}--\eqref{eq:ch2.reciprocal-impedance-matrix}
hold on \(\mathcal{K}_{\mathrm{acc}}\) after projection, while suppressed directions
Eq.~\eqref{eq:ch2.suppressed-angular-sheet-def} remain unobservable independently of reciprocity
\cite{stratton1941}. Angular budget
Eq.~\eqref{eq:ch2.angular-spectrum-composition-law} composes with port reciprocity without
eliminating null-space or geometric deficits \cite{stratton1941}.
Reciprocity acts on \(\operatorname{ran}\mathcal{R}_\omega\); geometric and observation
constraints from geometric constraints on radiative spectra (Subsection~\ref{sec:ch264}) and
Eq.~\eqref{eq:ch2.angular-spectrum-composition-law} restrict reported coordinates only
Laboratory errors and broken reciprocity. Magneto-electric media, active surfaces, and non-reciprocal port calibrations break
Eq.~\eqref{eq:ch2.transmit-receive-reciprocity} even when export maps differ
\(\mathcal{S}_\omega\) are accurate \cite{stratton1941}. Near-zone reactive
coupling measured without \(\Pi_{\mathrm{rad}}\) can appear non-reciprocal while For matched half-wave dipole ports at $k_0 a=1.2$ with $d=4a$, Eq.~\eqref{eq:ch2.reciprocal-impedance-matrix} requires $Z_{12}=Z_{21}$ when entries are extracted from $\mathcal{R}_\omega[\Jimp^{\mathrm{(tx)}}]$ on the far-zone branch. A typical reciprocity check gives $|Z_{12}-Z_{21}|/|Z_{11}|<0.02$ in reciprocal vacuum; swapping transmit and receive probes while holding $\mathcal{O}_{\mathrm{meas}}$ fixed reproduces the same mutual entry within numerical tolerance \cite{jackson1999}. If near-field probes at $k_0 r=1.5$ report $|Z_{12}/Z_{21}-1|\approx0.18$, the discrepancy is a regime failure, not a Lorentz violation: re-meter on $\mathcal{R}_\omega[\cdot]$ with $k_0 r>3$ before asserting broken reciprocity \cite{harrington2001}. Eq.~\eqref{eq:ch2.radiation-map-reciprocity} holds on \(\mathcal{F}_{\mathrm{rad}}\)
\cite{harrington2001}. Using different outgoing branches or terminations for
transmit and receive models violates the shared-operator assumption behind reciprocity in radiation and reception (Subsection~\ref{sec:ch283}) \cite{stratton1941}. Tikhonov inversion
Eq.~\eqref{eq:ch2.tikhonov-source-reconstruction} does not restore reciprocity where
\(\ker\mathcal{M}_\omega\) is nontrivial \cite{harrington2001}.
Subsection~\ref{sec:ch283} requires linear reciprocal media at fixed \(\omega\), passive
\(\mathcal{B}_\Gamma\), self-adjoint port functionals, and outgoing radiation condition
\cite{stratton1941}. It does not identify unique \(\Jimp\) from port data when
Eq.~\eqref{eq:ch2.reconstruction-null-space} is nontrivial \cite{stratton1941}. It does not
replace angular-spectrum operators in Section~\ref{sec:ch29} or harmonic bases in
\ref{ch:03} \cite{stratton1941}. Regime
qualification Eq.~\eqref{eq:ch2.regime-operator-screen} remains mandatory before transport
interpretation of port coupling \cite{stratton1941}. \medskip
\noindent
Section~\ref{sec:ch28} closes the reciprocity chain begun in its introduction:
Lorentz reciprocity in operator form (Subsection~\ref{sec:ch281}),
orthogonality and coupling of radiative modes (Subsection~\ref{sec:ch282}), and
reciprocity in radiation and reception (Subsection~\ref{sec:ch283})
\cite{stratton1941}. Spectral coordinates
from Section~\ref{sec:ch26} and null-space structure from Section~\ref{sec:ch27} now meet
exchange symmetry: bilinear coupling on \(\operatorname{ran}\mathcal{R}_\omega\) is constrained
by Eqs.~\eqref{eq:ch2.lorentz-reciprocity-operator}--\eqref{eq:ch2.reciprocal-impedance-matrix},
while \(\ker\mathcal{R}_\omega\), finite observation rank
Eq.~\eqref{eq:ch2.reconstruction-rank-cap}, and angular budget
Eq.~\eqref{eq:ch2.angular-spectrum-composition-law} explain deficits reciprocity alone does not
remove \cite{stratton1941}. Section~\ref{sec:ch29} develops plane-wave action and
retained/suppressed angular components on the same outgoing branch; Subsection~\ref{sec:ch294}
refines far-field angular accessibility; \ref{ch:03} supplies explicit harmonic
orthogonality used by reference through Subsection~3.4.3
\cite{harrington2001}. Transport and port-coupling claims must compose with
Eq.~\eqref{eq:ch2.section28-reciprocity-chain}, regime qualification
Eq.~\eqref{eq:ch2.regime-operator-screen}, and null-space chain
Eq.~\eqref{eq:ch2.section27-nullspace-chain}, not with raw \(\mathcal{S}_\omega[\Jimp]\) or
unqualified near-zone reactive content \cite{stratton1941}.
